\documentclass[12pt,a4paper]{article}

% --- Paquetes Básicos ---
\usepackage[utf8]{inputenc}
\usepackage[spanish]{babel}
\usepackage{amsmath}
\usepackage{amsfonts}
\usepackage{amssymb}
\usepackage{graphicx}
\usepackage{geometry}
\usepackage{listings}
\usepackage{xcolor}
\usepackage{hyperref}
\usepackage{booktabs}
\usepackage{caption}

% --- Configuración de Página ---
\geometry{margin=2.5cm}

% --- Estilo para Código ---
\definecolor{codegreen}{rgb}{0,0.6,0}
\definecolor{codegray}{rgb}{0.5,0.5,0.5}
\definecolor{codepurple}{rgb}{0.58,0,0.82}
\definecolor{backcolour}{rgb}{0.95,0.95,0.92}

\lstdefinestyle{mystyle}{
    backgroundcolor=\color{backcolour},
    commentstyle=\color{codegreen},
    keywordstyle=\color{magenta},
    numberstyle=\tiny\color{codegray},
    stringstyle=\color{codepurple},
    basicstyle=\ttfamily\footnotesize,
    breakatwhitespace=false,
    breaklines=true,
    captionpos=b,
    keepspaces=true,
    numbers=left,
    numbersep=5pt,
    showspaces=false,
    showstringspaces=false,
    showtabs=false,
    tabsize=2
}
\lstset{style=mystyle}

% --- Inicio del Documento ---
\begin{document}

\begin{titlepage}
    \centering
    \vspace*{1cm}
    \Huge
    \textbf{Análisis de Complejidad Computacional y Benchmarking Integral} \\
    \vspace{0.5cm}
    \LARGE
    Estudio Comparativo: Búsqueda y Ordenamiento \\
    \vspace{1.5cm}
    \textbf{Autor: [Tu Nombre]} \\
    \vfill
    \Large
    Ingeniería de Software \\
    \today
\end{titlepage}

\tableofcontents
\newpage

\begin{abstract}
Este informe documenta la evaluación empírica de algoritmos fundamentales de búsqueda y ordenamiento. El estudio abarca desde métodos de fuerza bruta ($O(n^2)$) hasta algoritmos de tipo "Divide y Vencerás" ($O(n \log n)$). A través de un análisis de regresión sobre tiempos de ejecución reales, se valida la jerarquía de eficiencia de la Notación Big O, demostrando por qué la elección algorítmica es crítica en el manejo de grandes volúmenes de datos.
\end{abstract}

\section{Marco Teórico}
La eficiencia asintótica permite predecir el comportamiento de un algoritmo conforme el tamaño de la entrada ($n$) tiende al infinito.

\subsection{Complejidades Evaluadas}
\begin{itemize}
    \item \textbf{Orden Logarítmico ($O(\log n)$):} Representado por la búsqueda binaria. Su crecimiento es extremadamente lento, lo que lo hace ideal para datasets masivos.
    \item \textbf{Orden Lineal ($O(n)$):} Representado por la búsqueda lineal. El tiempo de ejecución escala 1:1 con el tamaño de los datos.
    \item \textbf{Orden Cuasi-lineal ($O(n \log n)$):} Complejidad característica de \textit{QuickSort}. Es la eficiencia óptima para algoritmos de ordenamiento basados en comparaciones.
    \item \textbf{Orden Cuadrático ($O(n^2)$):} Métodos de \textit{Burbuja} y \textit{Selección}. El tiempo aumenta al cuadrado respecto a $n$, limitando su uso a datasets pequeños.
\end{itemize}

\section{Metodología de la Prueba}
Se implementó un entorno de pruebas en Python utilizando el generador de números aleatorios \texttt{random.randint} para crear datasets de diversos tamaños:
\begin{itemize}
    \item \textbf{Búsqueda:} Hasta $2 \times 10^6$ elementos.
    \item \textbf{Ordenamiento Eficiente:} Hasta $10^5$ elementos.
    \item \textbf{Ordenamiento Lento:} Hasta $5 \times 10^3$ elementos.
\end{itemize}

\section{Implementación de Algoritmos Relevantes}

\subsection{QuickSort: Divide y Vencerás}
El algoritmo QuickSort utiliza un pivote para particionar la lista en sub-listas de elementos menores y mayores, resolviendo el problema de forma recursiva.

\begin{lstlisting}[language=Python]
def quicksort(lista):
    if len(lista) <= 1: return lista
    pivote = lista[len(lista) // 2]
    izq = [x for x in lista if x < pivote]
    centro = [x for x in lista if x == pivote]
    der = [x for x in lista if x > pivote]
    return quicksort(izq) + centro + quicksort(der)
\end{lstlisting}

\section{Análisis de Resultados Empíricos}

% --- AQUÍ INSERTARÍAS TU GRÁFICA DE BÚSQUEDA ---
\subsection{Resultados de Búsqueda}
En las pruebas realizadas, la búsqueda binaria mantuvo tiempos cercanos a 0.0s incluso con 2 millones de datos. Por el contrario, la búsqueda lineal mostró una pendiente ascendente constante, sin diferencias significativas entre datos ordenados y desordenados en términos de complejidad.

% --- AQUÍ INSERTARÍAS TU GRÁFICA DE ORDENAMIENTO ---
\subsection{Resultados de Ordenamiento}
La comparativa entre \textit{Burbuja} y \textit{QuickSort} fue el hallazgo más notable. Mientras que \textit{Burbuja} superó los 3 segundos para solo 5,000 elementos, \textit{QuickSort} procesó 100,000 elementos en menos de 0.1 segundos, demostrando empíricamente la superioridad de $O(n \log n)$ sobre $O(n^2)$.

\section{Conclusiones}
\begin{enumerate}
    \item \textbf{Escalabilidad:} La arquitectura de software moderna debe priorizar algoritmos logarítmicos o cuasi-lineales para garantizar la escalabilidad.
    \item \textbf{Optimización:} El uso de \textit{QuickSort} reduce el tiempo de cómputo en varios órdenes de magnitud comparado con métodos tradicionales.
    \item \textbf{Hardware vs. Software:} A pesar de contar con hardware potente, un algoritmo ineficiente ($O(n^2)$) colapsará rápidamente frente al crecimiento de datos.
\end{enumerate}

\end{document}
